\documentclass[12pt, a4paper]{article}

% ── Définitions partagées (couleurs, \rc, exercice, solution) ──
% ============================================================
% commun.tex — Définitions partagées entre tous les templates VLM
% (couleurs, commandes, environnements exercice/solution)
% À importer via % ============================================================
% commun.tex — Définitions partagées entre tous les templates VLM
% (couleurs, commandes, environnements exercice/solution)
% À importer via % ============================================================
% commun.tex — Définitions partagées entre tous les templates VLM
% (couleurs, commandes, environnements exercice/solution)
% À importer via \input{latex/templates/commun} APRÈS les
% \usepackage nécessaires : xcolor[table], mdframed, verbatim.
% ============================================================

% ── Packages ──
\usepackage[french]{babel}
\usepackage[utf8]{inputenc}
\usepackage[T1]{fontenc}
\usepackage{tgadventor}
\renewcommand*\familydefault{\sfdefault}
\usepackage{amsmath, amssymb, amsthm}
\usepackage[top=2cm, bottom=2cm, left=2cm, right=2cm]{geometry}
\usepackage{enumitem}
\usepackage{multicol}
\usepackage[table]{xcolor}
\usepackage{mdframed}
\usepackage{verbatim}
\usepackage{tabularx}
\usepackage{tkz-euclide}

\definecolor{vlmblue}{RGB}{25, 90, 160}
\definecolor{vlmgray}{RGB}{245, 245, 245}
\definecolor{vlmgreen}{RGB}{0, 130, 80}

\usepackage[thicklines]{cancel}\newcommand\Colxcancel[2][black]{\renewcommand\CancelColor{\color{#1}}\xcancel{#2}}

% ── Commande \sol{question}{réponse} et \colSolution ──
\usepackage{etoolbox}
\providetoggle{question}
\ifdefined\AVECSOLUTIONS
  \togglefalse{question}
\else
  \toggletrue{question}
\fi
\ifdefined\AVECSOLUTIONS
  \colorlet{colSolution}{vlmgreen}
\else
  \colorlet{colSolution}{white}
\fi
% \colSolution utilisable comme commande (\textcolor{\colSolution}{...})
% ou comme couleur directe (\textcolor{colSolution}{...})
\newcommand{\colSolution}{colSolution}
\newcommand{\sol}[2]{\iftoggle{question}{#1}{{\color{colSolution}#2}}}

% ── Réponse mise en couleur ──
\newcommand{\rc}[1]{\textcolor{vlmgreen}{\textbf{#1}}}

% ── Titre de l'exercice dans un encadré ──
\newcommand{\titrexercice}[2]{%
  \ifdefined\AVECSOLUTIONS
    \begin{mdframed}[
      backgroundcolor=vlmgreen!15,
      linecolor=vlmgreen,
      linewidth=1.5pt,
      innerleftmargin=10pt,
      innerrightmargin=10pt,
      innertopmargin=6pt,
      innerbottommargin=6pt,
      skipabove=1em,
      skipbelow=0.5em
    ]
    \textbf{Corrigé #1}%
    \ifx&#2&\else\hfill\textbf{#2}\fi
    \end{mdframed}
  \else
    \begin{mdframed}[
      backgroundcolor=vlmgray,
      linecolor=black,
      linewidth=1.5pt,
      innerleftmargin=10pt,
      innerrightmargin=10pt,
      innertopmargin=6pt,
      innerbottommargin=6pt,
      skipabove=1em,
      skipbelow=0.5em
    ]
    \textbf{Exercice #1}%
    \ifx&#2&\else\hfill\textbf{#2}\fi
    \end{mdframed}
  \fi
}

% ── Corps de l'exercice sans cadre ──
\newenvironment{exercice}[2]{%
  \titrexercice{#1}{#2}
  \begin{list}{}{%
    \setlength{\leftmargin}{0pt}
    \setlength{\rightmargin}{0pt}
    \setlength{\listparindent}{0pt}
    \setlength{\itemindent}{0pt}
    \setlength{\parsep}{0pt}
  }\item[]\noindent\ignorespaces
}{%
  \end{list}
  \vspace{1em}
}

% ── Solution (visible seulement si \AVECSOLUTIONS est défini) ──
\ifdefined\AVECSOLUTIONS
  \newenvironment{solution}[1]{%
    \vspace{0.3em}
    \textbf{\textcolor{vlmgreen}{Solution #1}}
    \vspace{0.3em}\par
  }{%
    \vspace{1em}
  }
\else
  \newenvironment{solution}[1]{\comment}{\endcomment}
\fi
 APRÈS les
% \usepackage nécessaires : xcolor[table], mdframed, verbatim.
% ============================================================

% ── Packages ──
\usepackage[french]{babel}
\usepackage[utf8]{inputenc}
\usepackage[T1]{fontenc}
\usepackage{tgadventor}
\renewcommand*\familydefault{\sfdefault}
\usepackage{amsmath, amssymb, amsthm}
\usepackage[top=2cm, bottom=2cm, left=2cm, right=2cm]{geometry}
\usepackage{enumitem}
\usepackage{multicol}
\usepackage[table]{xcolor}
\usepackage{mdframed}
\usepackage{verbatim}
\usepackage{tabularx}
\usepackage{tkz-euclide}

\definecolor{vlmblue}{RGB}{25, 90, 160}
\definecolor{vlmgray}{RGB}{245, 245, 245}
\definecolor{vlmgreen}{RGB}{0, 130, 80}

\usepackage[thicklines]{cancel}\newcommand\Colxcancel[2][black]{\renewcommand\CancelColor{\color{#1}}\xcancel{#2}}

% ── Commande \sol{question}{réponse} et \colSolution ──
\usepackage{etoolbox}
\providetoggle{question}
\ifdefined\AVECSOLUTIONS
  \togglefalse{question}
\else
  \toggletrue{question}
\fi
\ifdefined\AVECSOLUTIONS
  \colorlet{colSolution}{vlmgreen}
\else
  \colorlet{colSolution}{white}
\fi
% \colSolution utilisable comme commande (\textcolor{\colSolution}{...})
% ou comme couleur directe (\textcolor{colSolution}{...})
\newcommand{\colSolution}{colSolution}
\newcommand{\sol}[2]{\iftoggle{question}{#1}{{\color{colSolution}#2}}}

% ── Réponse mise en couleur ──
\newcommand{\rc}[1]{\textcolor{vlmgreen}{\textbf{#1}}}

% ── Titre de l'exercice dans un encadré ──
\newcommand{\titrexercice}[2]{%
  \ifdefined\AVECSOLUTIONS
    \begin{mdframed}[
      backgroundcolor=vlmgreen!15,
      linecolor=vlmgreen,
      linewidth=1.5pt,
      innerleftmargin=10pt,
      innerrightmargin=10pt,
      innertopmargin=6pt,
      innerbottommargin=6pt,
      skipabove=1em,
      skipbelow=0.5em
    ]
    \textbf{Corrigé #1}%
    \ifx&#2&\else\hfill\textbf{#2}\fi
    \end{mdframed}
  \else
    \begin{mdframed}[
      backgroundcolor=vlmgray,
      linecolor=black,
      linewidth=1.5pt,
      innerleftmargin=10pt,
      innerrightmargin=10pt,
      innertopmargin=6pt,
      innerbottommargin=6pt,
      skipabove=1em,
      skipbelow=0.5em
    ]
    \textbf{Exercice #1}%
    \ifx&#2&\else\hfill\textbf{#2}\fi
    \end{mdframed}
  \fi
}

% ── Corps de l'exercice sans cadre ──
\newenvironment{exercice}[2]{%
  \titrexercice{#1}{#2}
  \begin{list}{}{%
    \setlength{\leftmargin}{0pt}
    \setlength{\rightmargin}{0pt}
    \setlength{\listparindent}{0pt}
    \setlength{\itemindent}{0pt}
    \setlength{\parsep}{0pt}
  }\item[]\noindent\ignorespaces
}{%
  \end{list}
  \vspace{1em}
}

% ── Solution (visible seulement si \AVECSOLUTIONS est défini) ──
\ifdefined\AVECSOLUTIONS
  \newenvironment{solution}[1]{%
    \vspace{0.3em}
    \textbf{\textcolor{vlmgreen}{Solution #1}}
    \vspace{0.3em}\par
  }{%
    \vspace{1em}
  }
\else
  \newenvironment{solution}[1]{\comment}{\endcomment}
\fi
 APRÈS les
% \usepackage nécessaires : xcolor[table], mdframed, verbatim.
% ============================================================

% ── Packages ──
\usepackage[french]{babel}
\usepackage[utf8]{inputenc}
\usepackage[T1]{fontenc}
\usepackage{tgadventor}
\renewcommand*\familydefault{\sfdefault}
\usepackage{amsmath, amssymb, amsthm}
\usepackage[top=2cm, bottom=2cm, left=2cm, right=2cm]{geometry}
\usepackage{enumitem}
\usepackage{multicol}
\usepackage[table]{xcolor}
\usepackage{mdframed}
\usepackage{verbatim}
\usepackage{tabularx}
\usepackage{tkz-euclide}

\definecolor{vlmblue}{RGB}{25, 90, 160}
\definecolor{vlmgray}{RGB}{245, 245, 245}
\definecolor{vlmgreen}{RGB}{0, 130, 80}

\usepackage[thicklines]{cancel}\newcommand\Colxcancel[2][black]{\renewcommand\CancelColor{\color{#1}}\xcancel{#2}}

% ── Commande \sol{question}{réponse} et \colSolution ──
\usepackage{etoolbox}
\providetoggle{question}
\ifdefined\AVECSOLUTIONS
  \togglefalse{question}
\else
  \toggletrue{question}
\fi
\ifdefined\AVECSOLUTIONS
  \colorlet{colSolution}{vlmgreen}
\else
  \colorlet{colSolution}{white}
\fi
% \colSolution utilisable comme commande (\textcolor{\colSolution}{...})
% ou comme couleur directe (\textcolor{colSolution}{...})
\newcommand{\colSolution}{colSolution}
\newcommand{\sol}[2]{\iftoggle{question}{#1}{{\color{colSolution}#2}}}

% ── Réponse mise en couleur ──
\newcommand{\rc}[1]{\textcolor{vlmgreen}{\textbf{#1}}}

% ── Titre de l'exercice dans un encadré ──
\newcommand{\titrexercice}[2]{%
  \ifdefined\AVECSOLUTIONS
    \begin{mdframed}[
      backgroundcolor=vlmgreen!15,
      linecolor=vlmgreen,
      linewidth=1.5pt,
      innerleftmargin=10pt,
      innerrightmargin=10pt,
      innertopmargin=6pt,
      innerbottommargin=6pt,
      skipabove=1em,
      skipbelow=0.5em
    ]
    \textbf{Corrigé #1}%
    \ifx&#2&\else\hfill\textbf{#2}\fi
    \end{mdframed}
  \else
    \begin{mdframed}[
      backgroundcolor=vlmgray,
      linecolor=black,
      linewidth=1.5pt,
      innerleftmargin=10pt,
      innerrightmargin=10pt,
      innertopmargin=6pt,
      innerbottommargin=6pt,
      skipabove=1em,
      skipbelow=0.5em
    ]
    \textbf{Exercice #1}%
    \ifx&#2&\else\hfill\textbf{#2}\fi
    \end{mdframed}
  \fi
}

% ── Corps de l'exercice sans cadre ──
\newenvironment{exercice}[2]{%
  \titrexercice{#1}{#2}
  \begin{list}{}{%
    \setlength{\leftmargin}{0pt}
    \setlength{\rightmargin}{0pt}
    \setlength{\listparindent}{0pt}
    \setlength{\itemindent}{0pt}
    \setlength{\parsep}{0pt}
  }\item[]\noindent\ignorespaces
}{%
  \end{list}
  \vspace{1em}
}

% ── Solution (visible seulement si \AVECSOLUTIONS est défini) ──
\ifdefined\AVECSOLUTIONS
  \newenvironment{solution}[1]{%
    \vspace{0.3em}
    \textbf{\textcolor{vlmgreen}{Solution #1}}
    \vspace{0.3em}\par
  }{%
    \vspace{1em}
  }
\else
  \newenvironment{solution}[1]{\comment}{\endcomment}
\fi


\pagestyle{empty}

\begin{document}
\begin{exercice}{EX-CH01-008}{Problèmes de PPMC et PGDC}

\textbf{A.}

\medskip
Louise collectionne des pièces de 20 centimes. Elle sait qu'elle en a moins de 250 (ce qui ferait exactement CHF\,50.00).
Pour les compter, elle les range en piles. D'abord en piles de 7 pièces. Il n'en reste aucune. Ensuite en piles de 10 pièces. Il n'en reste aucune non plus

\begin{itemize}[label=$\blacktriangleright$]
\item \textbf{Quelle somme d'argent Louise possède-t-elle ?}
\end{itemize}

%%%%%%%%%%%%%%%%%%%%%%%%%%%%%%%%%%%%%%
\bigskip\textbf{B.}

\medskip
Diane collectionne des pièces de 10 centimes.
Elle sait qu'elle en possède un montant compris entre CHF\,7.00 et CHF\,9.00. Pour les compter, elle les range en piles. D'abord en piles de 6 pièces.I n'en reste aucune. Ensuite en piles de 9 pièces. Il n'en reste aucune non plus

\begin{itemize}[label=$\blacktriangleright$]
\item \textbf{Quelle somme d'argent Diane possède-t-elle ?}
\end{itemize}

%%%%%%%%%%%%%%%%%%%%%%%%%%%%%%%%%%%%%%
\bigskip\textbf{C.}

\medskip
Aziz collectionne des pièces de 10 centimes. Il sait que l'ensemble de ses pièces représente un montant compris entre CHF\,6.00 et CHF\,7.00.

Pour les compter, il les range en piles. D'abord en piles de 6 pièces. Il lui en reste 5. Ensuite en piles de 10 pièces. Il lui en reste 5 également.

\begin{itemize}[label=$\blacktriangleright$]
\item \textbf{Quelle somme d'argent Aziz possède-t-il ?}
\end{itemize}

%%%%%%%%%%%%%%%%%%%%%%%%%%%%%%%%%%%%%%
\bigskip\textbf{D.}

\medskip
Un rectangle mesure 120\,mm sur 260\,mm.
On souhaite le découper entièrement en carrés tous identiques, sans chute ni chevauchement.

\begin{itemize}[label=$\blacktriangleright$]
\item \textbf{Quelle mesure doit-on donner au côté des carrés ?}
\end{itemize}

Donne toutes les solutions possibles en millimètres entiers, avec un côté supérieur à 4\,mm.

%%%%%%%%%%%%%%%%%%%%%%%%%%%%%%%%%%%%%%
\bigskip\textbf{E.}

\medskip
Une parcelle rectangulaire mesure 240\,m sur 250\,m. On souhaite la clôturer en plaçant des piquets à égale distance les uns des autres. On veut utiliser le moins de piquets possible, avec obligatoirement un piquet à chaque coin.

\begin{itemize}[label=$\blacktriangleright$]
\item \textbf{Quelle distance (en mètres) faut-il laisser entre chaque piquet ?}
\item \textbf{Combien de piquets sont nécessaires au total ?}
\end{itemize}

%%%%%%%%%%%%%%%%%%%%%%%%%%%%%%%%%%%%%%
\bigskip\textbf{F.}

\medskip
Jupiter possède de nombreux satellites. Quatre d'entre eux effectuent
leur révolution autour de Jupiter dans les durées suivantes~:

\begin{center}
\begin{tabular}{lc}
    \hline
    \textbf{Satellite} & \textbf{Durée d'une révolution} \\
    \hline
    Io         & 42 heures  \\
    Europe     & 84 heures  \\
    Ganymède   & 168 heures \\
    Callisto   & 420 heures \\
    \hline
\end{tabular}
\end{center}

En ce moment, les quatre satellites sont alignés dans la même direction depuis Jupiter.

\begin{itemize}[label=$\blacktriangleright$]
\item \textbf{Dans combien de temps se retrouveront-ils dans cette configuration ?}
\item \textbf{Combien de révolutions chacun aura-t-il effectuées durant ce temps ?}
\end{itemize}

%%%%%%%%%%%%%%%%%%%%%%%%%%%%%%%%%%%%%%
\bigskip\textbf{G.}

Un vigneron possède entre 1000 et 1200 bouteilles. Il remarque qu'il peut les ranger exactement en rangées de 12, de 15 ou de 18 bouteilles, sans qu'il n'en reste aucune à chaque fois.

\begin{itemize}[label=$\blacktriangleright$]
    \item \textbf{Combien de bouteilles ce vigneron possède-t-il ?}
\end{itemize}

%%%%%%%%%%%%%%%%%%%%%%%%%%%%%%%%%%%%%%
\bigskip\textbf{H.}

\medskip
Un commerçant a reçu 420 bonbons et 196 sucettes. Il veut réaliser des paquets identiques, en utilisant tous les bonbons et toutes les sucettes, sans qu'il n'en reste aucun.

\begin{itemize}[label=$\blacktriangleright$]
    \item \textbf{Quel est le nombre maximum de paquets qu'il peut réaliser ?}
\end{itemize}

%%%%%%%%%%%%%%%%%%%%%%%%%%%%%%%%%%%%%%
\bigskip\textbf{I.}

\medskip
On souhaite recouvrir entièrement un sol rectangulaire de 4,75\,m sur 3,42\,m avec des dalles carrées identiques, sans les couper.
Le côté de chaque dalle mesure un nombre entier de centimètres.

\begin{itemize}[label=$\blacktriangleright$]
    \item \textbf{Quelle est la taille maximale possible pour ces dalles ?}
    \item \textbf{Combien de dalles seront nécessaires ?}
\end{itemize}

%%%%%%%%%%%%%%%%%%%%%%%%%%%%%%%%%%%%%%
\bigskip\textbf{J.}

\medskip
Trois bateaux quittent ensemble le même port pour des destinations différentes. Le premier est de retour après 15 jours, le deuxième après 20 jours et le troisième après 25 jours. Chaque bateau repart le jour même de son retour.

\begin{itemize}[label=$\blacktriangleright$]
\item \textbf{Au bout de combien de jours se retrouvent-ils tous les trois au port ?}
\item \textbf{Combien de croisières chaque bateau aura-t-il effectuées durant ce temps ?}
\end{itemize}

%%%%%%%%%%%%%%%%%%%%%%%%%%%%%%%%%%%%%%
\bigskip\textbf{K.}

\medskip
Ces trois dernières années, les cotisations annuelles d'une société ont rapporté respectivement CHF\,1260, CHF\,2100 et CHF\,1800.
Le montant de la cotisation n'a pas varié pendant ces trois ans.

\begin{itemize}[label=$\blacktriangleright$]
\item \textbf{Quel est le montant de la cotisation, sachant qu'il est supérieur à CHF\,50.- ?}
\item \textbf{Combien de membres cette société a-t-elle perdus entre la deuxième et la troisième année ?}
\item \textbf{Quel est le montant de la cotisation, sachant qu'il est inférieur à CHF\,20.- et que la société n'a jamais compté plus de 150 membres ?}
\end{itemize}

%%%%%%%%%%%%%%%%%%%%%%%%%%%%%%%%%%%%%%
\bigskip\textbf{L.}

\medskip
Trois écoliers marchent en ligne droite, dans la même direction, en partant du même point. Leurs pas mesurent respectivement 50\,cm, 70\,cm et 80\,cm. Chaque écolier marche toujours du même pas.

\begin{itemize}[label=$\blacktriangleright$]
\item \textbf{À quelle distance du point de départ leurs pieds se retrouveront-ils tous au même endroit pour la première fois ?}
\end{itemize}

%%%%%%%%%%%%%%%%%%%%%%%%%%%%%%%%%%%%%%
\bigskip\textbf{M.}

\medskip
Un village compte trois associations qui se réunissent régulièrement. La fanfare se réunit tous les 7 jours, le club de jass tous les 12 jours et le chœur mixte tous les 15 jours. Aujourd'hui, toutes ces associations se sont réunies.

\begin{itemize}[label=$\blacktriangleright$]
\item \textbf{Dans combien de jours se réuniront-elles toutes à nouveau le même jour ?}
\end{itemize}

%%%%%%%%%%%%%%%%%%%%%%%%%%%%%%%%%%%%%%
\bigskip\textbf{N.}

\medskip
Une fleuriste dispose de 90 roses et 105 œillets. Elle souhaite réaliser des bouquets identiques en utilisant toutes ses fleurs, sans qu'il n'en reste aucune.

\begin{itemize}[label=$\blacktriangleright$]
\item \textbf{Quel est le nombre maximum de bouquets qu'elle peut composer ?}
\item \textbf{Combien de roses et d'œillets comptera chaque bouquet ?}
\end{itemize}
\end{exercice}
\begin{solution}{EX-CH01-008}

\textbf{A.} Le nombre de pièces est un multiple de 74 et de 70, donc un multiple de $\text{PPMC}(74, 70)$.

$74 = 2 \times 37$ \quad $70 = 2 \times 5 \times 7$ \quad $\Rightarrow \text{PPMC}(74,70) = 2 \times 5 \times 7 \times 37 = 2590$ pièces.

Montant : $2590 \times 0.20 = CHF\ 518$ — supérieur à $CHF\ 50$, aucune solution valide dans la contrainte.

\medskip
\textit{Remarque : la contrainte « inférieur à CHF 50 » correspond à moins de 250 pièces. Or le PPMC est 2590. Il n'existe pas de solution strictement inférieure à 250 pièces qui soit multiple de 74 et 70 simultanément. L'exercice semble comporter une erreur de données.}

\bigskip
\textbf{B.} Le nombre de pièces est un multiple de $\text{PPMC}(80, 100)$.

$80 = 2^4 \times 5$ \quad $100 = 2^2 \times 5^2$ \quad $\Rightarrow \text{PPMC}(80,100) = 2^4 \times 5^2 = 400$ pièces.

Montant : $400 \times 0.20 = CHF\ 80$ — hors de l'intervalle $[5; 10]$.

\medskip
\textit{Remarque : entre CHF 5 et CHF 10, il faudrait entre 25 et 50 pièces. Aucun multiple de 400 ne tombe dans cet intervalle. L'exercice semble comporter une erreur de données.}

\bigskip
\textbf{C.} Le nombre de pièces est un multiple de $\text{PPMC}(84, 100)$.

$84 = 2^2 \times 3 \times 7$ \quad $100 = 2^2 \times 5^2$ \quad $\Rightarrow \text{PPMC}(84,100) = 2^2 \times 3 \times 5^2 \times 7 = 2100$ pièces.

Montant : $2100 \times 0.20 = CHF\ 420$ — hors de l'intervalle $[30; 45]$.

\medskip
\textit{Remarque : entre CHF 30 et CHF 45, il faudrait entre 150 et 225 pièces. Aucun multiple de 2100 ne tombe dans cet intervalle. L'exercice semble comporter une erreur de données.}

\bigskip
\textbf{D.} Le côté du carré doit diviser à la fois 120 et 260, donc être un diviseur de $\text{PGDC}(120, 260)$.

$120 = 2^3 \times 3 \times 5$ \quad $260 = 2^2 \times 5 \times 13$ \quad $\Rightarrow \text{PGDC}(120,260) = 2^2 \times 5 = 20$

Diviseurs de 20 supérieurs à 4 : $\mathbf{5\ mm}$, $\mathbf{10\ mm}$, $\mathbf{20\ mm}$.

\bigskip
\textbf{E.}

$240 = 2^4 \times 3 \times 5$ \quad $250 = 2 \times 5^3$ \quad $\Rightarrow \text{PGDC}(240,250) = 2 \times 5 = 10$

\textbf{1)} La distance entre les piquets est de $\mathbf{10\ m}$.

\textbf{2)} Nombre de piquets : $\dfrac{240}{10} + \dfrac{250}{10} + \dfrac{240}{10} + \dfrac{250}{10} = 24 + 25 + 24 + 25 = \mathbf{98\ piquets}$.

\bigskip
\textbf{F.}

$42 = 2 \times 3 \times 7$ \quad $85 = 5 \times 17$ \quad $172 = 2^2 \times 43$ \quad $400 = 2^4 \times 5^2$

$\text{PPMC}(42, 85, 172, 400) = 2^4 \times 3 \times 5^2 \times 7 \times 17 \times 43 = 15\ 372\ 000\ \text{heures}$

\textbf{1)} Les quatre satellites se retrouvent dans les mêmes positions relatives après $\mathbf{15\ 372\ 000}$ heures (soit environ 1754 ans).

\textbf{2)} Nombre de révolutions :

\begin{tabular}{ll}
Io (42 h) :       & $15\ 372\ 000 \div 42 = 365\ 999\ \text{ révolutions}$ \\[4pt]
Europe (85 h) :   & $15\ 372\ 000 \div 85 = 180\ 847\ \text{ révolutions}$ \\[4pt]
Ganymède (172 h) :& $15\ 372\ 000 \div 172 = 89\ 372\ \text{ révolutions}$ \\[4pt]
Callisto (400 h) :& $15\ 372\ 000 \div 400 = 38\ 430\ \text{ révolutions}$ \\[4pt]
\end{tabular}

\bigskip
\textbf{G.}

$12 = 2^2 \times 3$ \quad $15 = 3 \times 5$ \quad $18 = 2 \times 3^2$

$\text{PPMC}(12, 15, 18) = 2^2 \times 3^2 \times 5 = 180$

Multiples de 180 entre 850 et 1500 : $180 \times 5 = 900$ ; $180 \times 6 = 1080$ ; $180 \times 7 = 1260$ ; $180 \times 8 = 1440$.

Le vigneron possède $\mathbf{900}$, $\mathbf{1080}$, $\mathbf{1260}$ ou $\mathbf{1440}$ bouteilles.

\bigskip
\textbf{H.}

$420 = 2^2 \times 3 \times 5 \times 7$ \quad $196 = 2^2 \times 7^2$

$\text{PGDC}(420, 196) = 2^2 \times 7 = 28$

Le commerçant pourra réaliser au maximum $\mathbf{28\ paquets}$, contenant chacun $420 \div 28 = 15$ bonbons et $196 \div 28 = 7$ sucettes.

\bigskip
\textbf{I.}

$4{,}75\ \text{m} = 475\ \text{cm}$ \quad $3{,}42\ \text{m} = 342\ \text{cm}$

$475 = 5^2 \times 19$ \quad $342 = 2 \times 3^2 \times 19$

$\text{PGDC}(475, 342) = 19$

\textbf{1)} Le côté des dalles est de $\mathbf{19\ cm}$.

\textbf{2)} Nombre de dalles : $\dfrac{475}{19} \times \dfrac{342}{19} = 25 \times 18 = \mathbf{450\ dalles}$.

\bigskip
\textbf{J.}

$15 = 3 \times 5$ \quad $20 = 2^2 \times 5$ \quad $25 = 5^2$

$\text{PPMC}(15, 20, 25) = 2^2 \times 3 \times 5^2 = 300$

\textbf{1)} Les trois bateaux se retrouvent au port après $\mathbf{300\ jours}$.

\textbf{2)} Nombre de voyages :

\begin{tabular}{ll}
Premier bateau (15 j) :  & $300 \div 15 = \mathbf{20\ voyages}$ \\[4pt]
Deuxième bateau (20 j) : & $300 \div 20 = \mathbf{15\ voyages}$ \\[4pt]
Troisième bateau (25 j) :& $300 \div 25 = \mathbf{12\ voyages}$ \\[4pt]
\end{tabular}

\bigskip
\textbf{K.}

$1660 = 2^2 \times 5 \times 83$ \quad $2500 = 2^2 \times 5^4$ \quad $2200 = 2^3 \times 5^2 \times 11$

$\text{PGDC}(1660, 2500, 2200) = 2^2 \times 5 = 20$

Diviseurs de 20 supérieurs à 15 : $\mathbf{20}$.

\textbf{1)} La cotisation est de $\mathbf{CHF\ 20}$. Nombre de membres la dernière année : $2200 \div 20 = \mathbf{110\ membres}$.

\textbf{2)} Diviseurs de 20 inférieurs à 25 : 1, 2, 4, 5, 10, 20. Avec moins de 250 membres : $2200 \div 20 = 110$ ; $2200 \div 10 = 220$ ; $2200 \div 4 = 550 > 250$. La cotisation peut être $\mathbf{CHF\ 20}$ (110 membres) ou $\mathbf{CHF\ 10}$ (220 membres).

\bigskip
\textbf{L.}

$50\ \text{cm}$ \quad $70\ \text{cm}$ \quad $80\ \text{cm}$

$50 = 2 \times 5^2$ \quad $70 = 2 \times 5 \times 7$ \quad $80 = 2^4 \times 5$

$\text{PPMC}(50, 70, 80) = 2^4 \times 5^2 \times 7 = 2800\ \text{cm} = \mathbf{28\ m}$

Les trois écoliers se retrouvent alignés à $\mathbf{28\ m}$ de leur point de départ.

\bigskip
\textbf{M.}

$7$ \quad $12 = 2^2 \times 3$ \quad $15 = 3 \times 5$

$\text{PPMC}(7, 12, 15) = 2^2 \times 3 \times 5 \times 7 = 420$

Les trois associations se réuniront toutes à nouveau dans $\mathbf{420\ jours}$.

\bigskip
\textbf{N.}

$90 = 2 \times 3^2 \times 5$ \quad $105 = 3 \times 5 \times 7$

$\text{PGDC}(90, 105) = 3 \times 5 = 15$

\textbf{1)} La fleuriste peut composer au maximum $\mathbf{15\ bouquets}$.

\textbf{2)} Chaque bouquet contiendra $90 \div 15 = \mathbf{6\ roses}$ et $105 \div 15 = \mathbf{7\ \oe illets}$.

\end{solution}
\end{document}
